\documentclass[11pt]{amsart}
\usepackage[margin=27mm]{geometry}
\usepackage{amsmath,amssymb,amsthm}
\usepackage{array,booktabs,longtable}
\usepackage[obeyspaces,spaces]{url}
\usepackage[colorlinks=true,linkcolor=blue,citecolor=blue,urlcolor=blue]{hyperref}
% The author's macros.
%
% some calligraphy
%
\newcommand\cA{{\mathcal A}}
\newcommand\cB{{\mathcal B}}
\newcommand\cC{{\mathcal C}}
\newcommand\cD{{\mathcal D}}
\newcommand\cE{{\mathcal E}}
\newcommand\cF{{\mathcal F}}
\newcommand\cG{{\mathcal G}}
\newcommand\cH{{\mathcal H}}
\newcommand\cI{{\mathcal I}}
\newcommand\cJ{{\mathcal J}}
\newcommand\cK{{\mathcal K}}
\newcommand\cL{{\mathcal L}}
\newcommand\cM{{\mathcal M}}
\newcommand\cN{{\mathcal N}}
\newcommand\cP{{\mathcal P}}
\newcommand\cQ{{\mathcal Q}}
\newcommand\cR{{\mathcal R}}
\newcommand\cS{{\mathcal S}}
\newcommand\cT{{\mathcal T}}
\newcommand\cU{{\mathcal U}}
\newcommand\cV{{\mathcal V}}

%
% fields and rings (and a semigroup)
%
\newcommand\cx{{\mathbb C}}% complexes
\newcommand\fld{{\mathbb F}}
\newcommand\flde{{\mathbb E}}
\newcommand\ints{{\mathbb Z}}
\newcommand\nn{{\mathbb N}}%non-negative integers
\newcommand\re{{\mathbb R}}%reals
\newcommand\rats{{\mathbb Q}}

\newcommand\Zu{{\mathbf u}}
\newcommand\Zv{{\mathbf v}}
\newcommand\Zw{{\mathbf w}}
\newcommand\Zy{{\mathbf y}}
\newcommand\Zz{{\mathbf z}}
\newcommand\Zx{{\mathbf x}}
\newcommand\Zj{{\mathbf j}}


%
% the really useful stuff
%
%\newcommand\opk[1]{\mathop{\hbox{\rm #1}}\nolimits} %% old version
\newcommand\opk[1]{\mathop{\mathrm{#1}}\nolimits}
\newcommand\comp[1]{{\mkern2mu\overline{\mkern-2mu#1}}}
\newcommand\diff{\mathbin{\mkern-1.5mu\setminus\mkern-1.5mu}}% for \setminus
\newcommand\gbin[2]{\genfrac{[}{]}{0pt}{}{#1}{#2}}
\newcommand\inv{^{-1}}
\newcommand\res{\mathbin{\mkern-2.0mu\restriction\mkern-2.0mu}}% ch2
\newcommand\sbs{\subseteq}
\newcommand\sps{\supseteq}
\newcommand\seq[3]{#1_{#2},\ldots,#1_{#3}}
\newcommand\supp{\opk{supp}}
\newcommand\im{\opk{im}}
\newcommand\row{\opk{row}}
\newcommand\pmat[1]{\begin{pmatrix} #1 \end{pmatrix}}
\newcommand\cprod{\mathbin{\square}}
%
% matrix theory
%
\newcommand\ip[2]{\langle#1,#2\rangle}
\newcommand\one{{\bf1}}
\newcommand\rk{\opk{rk}}
\newcommand\tr{\opk{tr}}
\newcommand\col{\opk{col}}
\newcommand\mat[3]{\mathrm{Mat}_{#1\times #2}(#3)}

%
% some group theory
%
\newcommand\aut[1]{{\rm Aut}(#1)}
\newcommand\fx[1]{{\rm fix}(#1)}% ch2
\newcommand\grp[1]{\langle #1\rangle}
\newcommand\nrml{\vartriangleleft}
\newcommand\nrmleq{\trianglelefteq}
\newcommand\alt[1]{\opk{Alt}(#1)}
\DeclareMathOperator{\Sym}{Sym}
% End of the author's macros.
\newtheorem{theorem}{Theorem}[section]
\newtheorem{lemma}[theorem]{Lemma}
\newtheorem{proposition}[theorem]{Proposition}
\newtheorem{corollary}[theorem]{Corollary}
\theoremstyle{definition}
\newtheorem{definition}[theorem]{Definition}
\newtheorem{remark}[theorem]{Remark}
\newcommand{\typ}[1]{\mathrm{#1}}
\newcommand{\Zr}{{\mathbf r}}
\newcommand{\Zs}{{\mathbf s}}
\newcommand{\Zt}{{\mathbf t}}
\newcommand{\Ze}{{\mathbf e}}
\newcommand{\Zb}{{\mathbf b}}
\DeclareRobustCommand{\ml}[1]{\ifmmode\text{\normalfont\small\ttfamily\detokenize{#1}}\else{\normalfont\small\nolinkurl{#1}}\fi}
\emergencystretch=2em
\title[An ordinary proof of Theorem 3.12]{An ordinary proof of Theorem 3.12 of Cameron, Goethals, Seidel and Shult}
\date{28 September 2026}

\begin{document}
\raggedbottom
\begin{abstract}
We characterize the sets of lines at angles $90^\circ$, $60^\circ$ and
$45^\circ$ which are contained in the lines of a finite reduced
crystallographic Root System. A suitable partition assigns squared lengths
2 and 4 to the initial vectors. In their integer span, we take all vectors
of squared length 2 and those of squared length 4 which pair evenly with
the span. We prove that this finite set satisfies all Root System axioms.
Conversely, crystallographic integrality gives a partition by the parity
of the integer logarithm of squared length. We include the empty
and nonspanning cases, exact computational checks, and the statement shape
of the Lean formalization.
\end{abstract}
\maketitle

\section{Statement and conventions}\label{sec:statement}

We will write $V$ for a finite-dimensional real inner product space and
$(\Zx,\Zy)$ for its inner product. Set $q(\Zx)=(\Zx,\Zx)$.
A \emph{line} is a one-dimensional linear subspace of $V$; in particular,
every line passes through zero. For nonzero $\Zx\in\ell$ and
$\Zy\in m$, set
\begin{equation}\label{eq:cosine}
 c(\ell,m)=\frac{(\Zx,\Zy)^2}{q(\Zx)q(\Zy)}.
\end{equation}
Replacing either representative by a nonzero multiple multiplies the
numerator and denominator by the same square. Thus $c$ is well-defined.
For distinct lines, the angles $90^\circ$, $60^\circ$ and $45^\circ$
mean, respectively, $c=0$, $c=1/4$ and $c=1/2$.
These are the unoriented line angles in $[0,90^\circ]$. We will use their squared
cosines throughout, so the argument does not require trigonometry.
Entries of matrices and vectors are function values, such as $M(i,j)$
and $\Zx(i)$; subscripts index families of vectors.

\begin{definition}\label{def:rootset}
\textup{\cite[Definition 3.9, with the corrections below]{CGSS}.}
A \emph{Root System} in $V$ is a finite set
$\Phi\sbs V\diff\{0\}$ which spans $V$ over $\re$ and satisfies:
\begin{enumerate}
\item for every $\Zr,\Zs\in\Phi$,
      $2(\Zr,\Zs)/q(\Zr)\in\ints$;
\item if $\Zr\in\Phi$, $t\in\re$ and $t\Zr\in\Phi$, then
      $t\in\{1,-1\}$; and
\item for every $\Zr,\Zs\in\Phi$,
\[
 \tau_{\Zr}(\Zs):=\Zs-\frac{2(\Zr,\Zs)}{q(\Zr)}\Zr\in\Phi.
\]
\end{enumerate}
We will denote its set of lines by
$\cL(\Phi)=\{\re\Zr:\Zr\in\Phi\}$.
\end{definition}

This is the finite, reduced, crystallographic convention, with all three
qualifiers required. The nonzero condition makes every denominator
positive. In the printed condition (b), we read $S$ as $\Sigma$.
The same definition is used in \cite[Theorems A--C]{R311}.
We allow $\Phi=\varnothing$ when $V=0$. Reflection in a root sends
that root to its negative, so every root line has both signs; reducedness
then says that these are its only two roots.

\begin{theorem}[Theorem 3.12 of \cite{CGSS}, made precise]\label{thm:main}
Let $V$ be a finite-dimensional real inner product space, and let $S$ be
a set of one-dimensional linear subspaces of $V$. Suppose that for every
pair of distinct lines $\ell,m\in S$,
\[
 c(\ell,m)\in\{0,1/4,1/2\}.
\]
The following conditions are equivalent.
\begin{enumerate}
\item There is a Root System $\Phi$ spanning $V$ such that
      $S\sbs\cL(\Phi)$.
\item There are disjoint subsets $\cP,\cQ$ with
      $S=\cP\cup\cQ$ such that distinct lines in the same part have
      $c\in\{0,1/4\}$, and lines in different parts have
      $c\in\{0,1/2\}$.
\end{enumerate}
\end{theorem}

\begin{remark}\label{rem:conventions}
Lemma~\ref{lem:finite-lines} proves that $S$ is automatically finite.
The set $S$ and either part may be empty, and $S$ need not span $V$.
Allowing an empty part is necessary: a single line lies in a Root System
but has no partition into two nonempty parts, and two $60^\circ$ lines
in $\typ A_2$ must lie in the same part.
This is the proposed Implicit hypothesis 8 in \ml{cgss/errata.md}.
Allowing a containing Root System in a larger inner product space gives
the same condition: the necessary-direction proof uses only the inner
products of the selected roots, and the sufficient direction constructs
a Root System in $V$ itself.
\end{remark}

In the last clause of the
printed theorem, ``$90^\circ$ and $45^\circ$'' means
``$90^\circ$ or $45^\circ$''; this is Error 13, item 3 of the project
errata. We prove the equivalence missing in Gap 11.
Corollary~\ref{cor:normalized-containment} records the stronger construction
with squared root lengths in $\{2,4\}$, both in the span and in $V$.

Our proof follows route RS3 in \ml{cgss/prerequisites.md}, Sections 6.2
and 13. We will prove both directions directly. The classification in
\cite[Theorems A and B]{R311} is an allowed prerequisite, but is not
needed for this route. The root-angle calculation in its Theorem C is
the one we will write out in Lemma~\ref{lem:root-angles}. We will also
give the integer-span finiteness argument referred to as LN1 in the
prerequisites. Thus we use no unproved completion or root-system
classification assertion from \cite{CGSS}.

\section{Finiteness and the integer span}\label{sec:finite}

\begin{lemma}\label{lem:finite-lines}
Let $V$ be a finite-dimensional real inner product space, and let $S$ be
a set of one-dimensional linear subspaces whose distinct-line squared
cosines belong to $\{0,1/4,1/2\}$. Then $S$ is finite.
\end{lemma}
\begin{proof}
Let $U$ be the set of all unit vectors on the lines of $S$, including
both signs, and let $W$ be its real span. Every line has a unit vector:
divide any nonzero vector on it by its positive norm. Thus $W$ is also
the span of all the lines. Finite-dimensional linear algebra supplies
a basis $\Zb_1,\ldots,\Zb_d$ of $W$ chosen from $U$.
If $U$ is empty, use the empty basis and $W=0$.
For each $\Zx\in U$ and each $i$, either $\Zx$ is on the line of
$\Zb_i$, when $\Zx=\pm\Zb_i$, or the two lines are distinct. In
the second case \eqref{eq:cosine} applies. Thus
\[
 (\Zx,\Zb_i)\in\Theta,
 \qquad\Theta=\{-1,-1/\sqrt2,-1/2,0,1/2,1/\sqrt2,1\}.
\]
The map $\Zx\mapsto((\Zx,\Zb_i))_{i=1}^d$ is injective on $W$.
Indeed, if $\Zx,\Zy$ have the same image, their difference is
orthogonal to every basis vector. By bilinearity it is orthogonal to all
of $W$. Since $\Zx-\Zy\in W$, its squared norm is zero,
so $\Zx=\Zy$. This embeds $U$ into the finite set $\Theta^d$.
The map $\Zx\mapsto\re\Zx$ maps $U$ onto $S$, proving finiteness.
This argument also covers $d=0$, when $W$ has no unit vector.
\end{proof}

\begin{lemma}\label{lem:integer-span}
Let $T=\{\Zt_1,\ldots,\Zt_m\}$ be a finite set of vectors in a
real inner product space. Suppose every $(\Zt_i,\Zt_j)$ is an
integer and every $q(\Zt_i)$ belongs to $\{2,4\}$. Put
\[
 L=\left\{\sum_{i=1}^m a(i)\Zt_i:a(i)\in\ints\right\},
 \qquad W=\opk{span}_{\re}(T).
\]
Then every inner product of elements of $L$ is an integer, and
$\{\Zx\in L:q(\Zx)\leq4\}$ is finite. More precisely, its
cardinality is at most $9^d$, where $d=\dim W$.
If $(\Zt_i,\Zt_j)$ is even for every $j$, then
$(\Zt_i,\Zx)$ is even for every $\Zx\in L$.
\end{lemma}
\begin{proof}
For integer coefficient functions $a,b$, bilinearity gives
\[
 \left(\sum_i a(i)\Zt_i,\sum_j b(j)\Zt_j\right)
 =\sum_{i,j}a(i)b(j)(\Zt_i,\Zt_j)\in\ints.
\]
The same calculation with $\Zt_i$ fixed proves the last assertion.
No independence of $T$ is needed: any expressions for the vectors
give the required integer or even integer.

Choose a basis $\Zb_1,\ldots,\Zb_d$ of $W$ from $T$.
For $\Zx\in L$ with $q(\Zx)\leq4$, Cauchy--Schwarz gives
\[
 (\Zx,\Zb_i)^2\leq q(\Zx)q(\Zb_i)\leq16.
\]
Thus each $(\Zx,\Zb_i)$ is an integer in $\{-4,\ldots,4\}$.
The inner-product coordinate map is injective by the argument in
Lemma~\ref{lem:finite-lines}. Its target has $9^d$ elements.
If $T$ is empty, then $L=W=\{0\}$, and the assertion gives the
correct bound $9^0=1$.
\end{proof}

The integral inner products are essential to this argument. We do not
infer discreteness merely from finite generation of an additive subgroup
of a real vector space. The coordinate proof is the bounded-norm form
of LN1, and supplies its proof here.

\section{An explicit Root System for the sufficient direction}\label{sec:construction}

\begin{lemma}\label{lem:reflection}
Let $V$ be a real inner product space and $\Zr\in V$ be nonzero. Then
the map $\tau_{\Zr}(\Zx)=\Zx-2(\Zr,\Zx)\Zr/q(\Zr)$ is linear,
is an involution, preserves squared norms, and satisfies
\begin{equation}\label{eq:self-adjoint}
 (\tau_{\Zr}\Zx,\Zy)=(\Zx,\tau_{\Zr}\Zy)
 \qquad(\Zx,\Zy\in V).
\end{equation}
Also $\tau_{\Zr}\Zr=-\Zr$.
\end{lemma}
\begin{proof}
Linearity follows from bilinearity of the inner product. Direct expansion
gives
\begin{align*}
 (\Zr,\tau_{\Zr}\Zx)&=-(\Zr,\Zx),\\
 \tau_{\Zr}(\tau_{\Zr}\Zx)
 &=\Zx-\frac{2(\Zr,\Zx)}{q(\Zr)}\Zr
       +\frac{2(\Zr,\Zx)}{q(\Zr)}\Zr=\Zx,\\
 q(\tau_{\Zr}\Zx)
 &=q(\Zx)-\frac{4(\Zr,\Zx)^2}{q(\Zr)}
              +\frac{4(\Zr,\Zx)^2}{q(\Zr)}=q(\Zx),\\
 (\tau_{\Zr}\Zx,\Zy)
 &=(\Zx,\Zy)-\frac{2(\Zr,\Zx)(\Zr,\Zy)}{q(\Zr)}
  =(\Zx,\tau_{\Zr}\Zy).
\end{align*}
The denominator is nonzero because $\Zr\ne0$.
Substitution of $\Zx=\Zr$ gives $\tau_{\Zr}\Zr=-\Zr$.
\end{proof}

\begin{proposition}\label{prop:psi}
Let $V$ be a real inner product space, and let $T$ be a finite set of
vectors such that:
\begin{enumerate}
\item $q(\Zt)\in\{2,4\}$ for every $\Zt\in T$;
\item $(\Zt,\Zu)\in\ints$ for every $\Zt,\Zu\in T$; and
\item if $\Zt\in T$ has $q(\Zt)=4$, then
      $(\Zt,\Zu)\in2\ints$ for every $\Zu\in T$.
\end{enumerate}
Let $L=\opk{span}_{\ints}(T)$ and $W=\opk{span}_{\re}(T)$.
The set
\begin{equation}\label{eq:psi}
\begin{split}
 \Psi(T)={}&\{\Zx\in L:q(\Zx)=2\}\\
 &{}\cup\{\Zx\in L:q(\Zx)=4\text{ and }
                 (\Zx,\Zy)\in2\ints\text{ for every }\Zy\in L\}
\end{split}
\end{equation}
is a Root System in $W$ containing $T$, with squared root lengths in
$\{2,4\}$.
\end{proposition}
\begin{proof}
Write $\Psi=\Psi(T)$. Hypotheses (1) and (3), together with
Lemma~\ref{lem:integer-span}, give $T\sbs\Psi\sbs L\sbs W$.
Thus $\Psi$ spans $W$. Its vectors are nonzero, and it is a subset
of the finite set $\{\Zx\in L:q(\Zx)\leq4\}$ in that lemma.

Take $\Zr\in\Psi$. For every $\Zx\in L$, the reflection coefficient is
\begin{equation}\label{eq:coefficient}
 \frac{2(\Zr,\Zx)}{q(\Zr)}
 =\begin{cases}
   (\Zr,\Zx),&q(\Zr)=2;\\
   (\Zr,\Zx)/2,&q(\Zr)=4.
  \end{cases}
\end{equation}
In the first case it is an integer by Lemma~\ref{lem:integer-span}.
In the second case it is an integer by the evenness condition in
\eqref{eq:psi}.
Since $L$ is closed under subtraction and integer multiples,
$\tau_{\Zr}(\Zx)\in L$. The involution property in
Lemma~\ref{lem:reflection} then gives $\tau_{\Zr}(L)=L$.

Take $\Zs\in\Psi$ and set $\Zt=\tau_{\Zr}(\Zs)$.
We have just proved $\Zt\in L$. By norm preservation,
$q(\Zt)=q(\Zs)\in\{2,4\}$. If $q(\Zt)=4$, then
$q(\Zs)=4$; for every $\Zx\in L$, \eqref{eq:self-adjoint} gives
\[
 (\Zt,\Zx)=(\Zs,\tau_{\Zr}\Zx)\in2\ints,
\]
since $\tau_{\Zr}\Zx\in L$. Thus $\Zt\in\Psi$ in either length
case, proving reflection closure.

For $\Zr,\Zs\in\Psi$, formula \eqref{eq:coefficient}, applied
to $\Zs\in L$, proves crystallographic integrality. We will check
reducedness explicitly. Suppose $\Zs=t\Zr$ with
$\Zr,\Zs\in\Psi$. Both vectors are nonzero. If their squared
lengths are equal, $q(\Zs)=t^2q(\Zr)$ gives $t^2=1$, so
$t=1$ or $t=-1$. If their squared lengths differ, relabel them
and replace $t$ by $1/t$ if needed so that $q(\Zr)=2$ and
$q(\Zs)=4$. Their proportionality then
gives
\[
 t^2=2,\qquad t=\frac{(\Zr,\Zs)}2\in\rats.
\]
The first equation implies $t=\sqrt2$ or $t=-\sqrt2$, contrary
to the irrationality of $\sqrt2$. Thus different lengths cannot
occur on one line. This proves reducedness.

We have now proved finiteness, nonzero roots, spanning, integrality,
reducedness and reflection closure, which are exactly
Definition~\ref{def:rootset} in $W$. If $T$ is empty, then $L=\{0\}$,
$\Psi=\varnothing$ and $W=0$, giving the permitted empty Root System.
\end{proof}

\begin{lemma}\label{lem:seed-completion}
Let $V$ be a finite-dimensional real inner product space and let
$T\sbs V$ be finite. Suppose that every vector in $T$ has squared
length 2 or 4, that all inner products in $T$ are integers, and that
each vector of squared length 4 pairs evenly with every vector in $T$.
Put $W=\opk{span}_{\re}(T)$, and let $\Ze_1,\ldots,\Ze_h$ be an
orthonormal basis of $W^\perp$. Then
\[
 T'=T\cup\{\sqrt2\Ze_i:1\leq i\leq h\}
\]
is finite, satisfies the same three conditions, and spans $V$.
\end{lemma}
\begin{proof}
Each new vector has squared length 2. Its inner products with $T$
and with distinct new vectors are zero; its inner product with itself
is 2. Thus all new inner products are
integers, and the new pairings involving an old vector of squared length
4 are even. No new vector has squared length 4. This proves the three
conditions. The finite union $T'$ spans $W+W^\perp=V$ by orthogonal
decomposition and $\sqrt2\ne0$. If $W=V$, the new basis is empty;
if $W=0$, it is a basis of $V$. This also includes $V=0$.
\end{proof}

\begin{corollary}\label{cor:normalized-containment}
Let $V$ be a finite-dimensional real inner product space and let
$S=\cP\cup\cQ$ be a set of lines, with disjoint parts. Suppose that
distinct lines in the same part have squared cosine 0 or $1/4$, and
lines in different parts have squared cosine 0 or $1/2$.
There is a Root System in the span of $S$ whose lines contain $S$
and whose squared root lengths belong to $\{2,4\}$. There is also
such a Root System spanning $V$.
\end{corollary}
\begin{proof}
Lemma~\ref{lem:finite-lines} makes $S$ finite. Choose one vector
$\Zt_\ell$ on each line, of squared length 2 if $\ell\in\cP$
and squared length 4 if $\ell\in\cQ$. Such a vector is obtained
by normalizing any nonzero vector on the line and multiplying by
$\sqrt2$ or 2. Either sign is permitted. Put
$T=\{\Zt_\ell:\ell\in S\}$.
For distinct lines, \eqref{eq:cosine} and the partition conditions give
\[
\begin{array}{c|c|c|c}
\hbox{parts}&q(\Zt_\ell)q(\Zt_m)&c(\ell,m)&(\Zt_\ell,\Zt_m)\\\hline
\cP,\cP&4&0\hbox{ or }1/4&0,\ 1,\ -1\\
\cQ,\cQ&16&0\hbox{ or }1/4&0,\ 2,\ -2\\
\cP,\cQ&8&0\hbox{ or }1/2&0,\ 2,\ -2.
\end{array}
\]
On the diagonal the inner product is 2 or 4. Thus every inner product
is an integer, and every inner product involving a vector of squared
length 4 is even. Proposition~\ref{prop:psi} gives the Root System
$\Psi(T)$ in the span $W$ of $S$. Each chosen vector lies in
$\Psi(T)$, so its lines contain $S$.
For the ambient assertion, take an orthonormal basis of $W^\perp$ and
form $T'$ as in Lemma~\ref{lem:seed-completion}. Apply
Proposition~\ref{prop:psi} directly to $T'$ in $V$. Its output
$\Psi(T')$ spans $V$ and contains $T$, so its lines also contain $S$.
Both applications give squared root lengths in $\{2,4\}$.
\end{proof}

\begin{proof}[Proof of Theorem~\ref{thm:main}, (2) implies (1)]
The ambient assertion of Corollary~\ref{cor:normalized-containment}
gives condition (1).
\end{proof}

\section{Root lengths and the necessary direction}\label{sec:necessity}

\begin{lemma}\label{lem:root-angles}
\textup{\cite[Theorem C, root-angle assertion]{R311}, in the two cases used here.}
Let $\Phi$ be a Root System in a real inner product space, and let
$\Zr,\Zs\in\Phi$ span distinct lines. If their line angle is
$60^\circ$, then $q(\Zs)=q(\Zr)$. If it is $45^\circ$, then
$q(\Zs)/q(\Zr)\in\{2,1/2\}$.
\end{lemma}
\begin{proof}
Put $a=2(\Zr,\Zs)/q(\Zr)$ and
$b=2(\Zr,\Zs)/q(\Zs)$. By crystallographic integrality both are
integers. For either angle they are nonzero and have the same sign.
Moreover,
\[
 ab=\frac{4(\Zr,\Zs)^2}{q(\Zr)q(\Zs)}
    =4c(\re\Zr,\re\Zs),
 \qquad \frac{a}{b}=\frac{q(\Zs)}{q(\Zr)}.
\]
At $60^\circ$, $ab=1$; the positive integers $|a|,|b|$ must both
be 1, so $a/b=1$. At $45^\circ$, $ab=2$, so their ordered absolute
values are $(1,2)$ or $(2,1)$, giving $a/b=1/2$ or 2.
This proves the assertions directly from Definition~\ref{def:rootset}.
\end{proof}

\begin{proof}[Proof of Theorem~\ref{thm:main}, (1) implies (2)]
Let $\Phi$ be a Root System in $V$ whose lines contain $S$.
For each $\ell\in S$, choose a root $\Zr_\ell\in\Phi$ on it and
put $\delta(\ell)=q(\Zr_\ell)>0$. Reducedness and the presence of both
signs show that $\delta(\ell)$ is independent of the chosen root.
For $x>0$, let $\eta(x)$ be the unique integer with
\begin{equation}\label{eq:log-bounds}
 2^{\eta(x)}\leq x<2^{\eta(x)+1}.
\end{equation}
This is $\lfloor\log_2 x\rfloor$, implemented by \ml{Int.log 2 x}.
The integer-logarithm bounds in Section~\ref{sec:sources} supply
existence; strict increase of $j\mapsto2^j$ gives uniqueness.
Multiplying \eqref{eq:log-bounds} by 2 gives
$2^{\eta(x)+1}\leq2x<2^{\eta(x)+2}$. Uniqueness, followed by the same
identity with $x/2$ in place of $x$, gives
\begin{equation}\label{eq:log-shift}
 \eta(2x)=\eta(x)+1,\qquad \eta(x/2)=\eta(x)-1.
\end{equation}

Put in $\cP$ the lines for which $\eta(\delta(\ell))$ is even, and
in $\cQ$ those for which it is odd. Every integer is exactly one of
even and odd, including negative integers, so these parts are disjoint
and have union $S$. By Lemma~\ref{lem:root-angles}, a $60^\circ$
pair has equal squared lengths and equal parities. A $45^\circ$ pair
has squared lengths in ratio 2 or $1/2$, so \eqref{eq:log-shift}
makes its parities opposite. Orthogonal pairs impose no condition.
Thus two lines in the same part have angle $90^\circ$ or $60^\circ$,
and two lines in different parts have angle $90^\circ$ or $45^\circ$.
For $S=\varnothing$ take both parts empty. This proves condition (2)
and completes the proof of Theorem~\ref{thm:main}.
\end{proof}

The colouring is defined globally, with no choice of components or paths.
Unrelated scales on orthogonal components impose no additional constraint,
so a decomposable Root System may have more than two lengths globally.
A containing component of type $\typ G_2$ is also permitted: the angle
restriction applies only to the selected lines.

\section{Examples and computational checks}\label{sec:checks}

\subsection{The role of the parity invariant}

We will first give a small example showing why \eqref{eq:psi} restricts
the vectors of squared length 4. Take independent vectors
$\Zb_1,\Zb_2,\Zb_3$ with Gram matrix
\[
 M=\pmat{2&-1&0\\-1&2&0\\0&0&2}.
\]
Such vectors exist: for an orthonormal basis $\Ze_1,\Ze_2,\Ze_3$,
take $\Zb_1=\sqrt2\Ze_1$,
$\Zb_2=-\Ze_1/\sqrt2+\sqrt{3/2}\Ze_2$ and
$\Zb_3=\sqrt2\Ze_3$.
All initial vectors have squared length 2 and have integral inner
products. In their integer span the vector $\Zr=\Zb_1+\Zb_3$
has squared length 4, but $(\Zr,\Zb_2)=-1$. Thus
\[
 \frac{2(\Zr,\Zb_2)}{q(\Zr)}=-\frac12,
 \qquad
 \tau_{\Zr}(\Zb_2)=\Zb_2+\frac12(\Zb_1+\Zb_3)\notin L.
\]
The nonmembership follows from independence of the basis and the
noninteger coefficients. Thus taking every vector of squared length
2 or 4 in $L$ would fail. The evenness condition excludes this $\Zr$
from $\Psi(T)$.

The set $\Psi(T)$ need not be the least reflection-closed set containing
$T$. For two orthogonal seeds of squared length 2, it has eight roots,
whereas reflection closure has four. For the usual $\typ D_4$ seeds,
$\Psi(T)$ has 48 roots of type $\typ F_4$, whereas reflection closure
has 24. We will retain reflection closure as a separate computational
comparison below; its counts are not counts of $\Psi(T)$.

\subsection{Classical models and the exceptional model}

We will use the coordinate models proved in \cite[Theorem B and the
coordinate-model section]{R311}. Let $\Ze_1,\ldots,\Ze_k$ be
the standard orthonormal basis. Their root sets are
\begin{align*}
 \Phi(\typ B_k)
  &=\{\pm\Ze_i:1\leq i\leq k\}
       \cup\{\pm\Ze_i\pm\Ze_j:i<j\},\\
 \Phi(\typ C_k)
  &=\{\pm2\Ze_i:1\leq i\leq k\}
       \cup\{\pm\Ze_i\pm\Ze_j:i<j\},\\
 \Phi(\typ F_4)
  &=\{\pm\Ze_i:1\leq i\leq4\}
       \cup\{\pm\Ze_i\pm\Ze_j:i<j\}\\
  &\hspace{12mm}\cup
       \left\{\frac12\sum_{i=1}^4\epsilon(i)\Ze_i:
                   \epsilon(i)\in\{1,-1\}\right\}.
\end{align*}
In $\typ B_k$, short roots are on the axes; in $\typ C_k$,
long roots are on the axes. The two line sets coincide, while the
length partitions are interchanged. Multiplying all roots of
$\typ B_k$ or $\typ F_4$ by $\sqrt2$ gives squared lengths
2 and 4. The $\typ C_k$ model already has that normalization.
The allowed squared cosines and the length partitions satisfy
Theorem~\ref{thm:main}, by the root-angle calculation or directly
from the coordinate lists.

The script \ml{checks/check_thm312.py} uses the following generating
seeds before that common normalization:
\[
\begin{array}{c|l}
\typ B_k&\Ze_1-\Ze_2,\ldots,\Ze_{k-1}-\Ze_k,\Ze_k\\
\typ C_k&\Ze_1-\Ze_2,\ldots,\Ze_{k-1}-\Ze_k,2\Ze_k\\
\typ F_4&\Ze_2-\Ze_3,\Ze_3-\Ze_4,\Ze_4,
             (\Ze_1-\Ze_2-\Ze_3-\Ze_4)/2.
\end{array}
\]
It verifies all Root System axioms on the full coordinate sets,
computes the partition from line angles, and checks its agreement with
the two length classes. Independently, it constructs reflection closure
from the displayed seeds with the normalized integer Gram matrix,
checking membership in $\Psi(T)$ at every new root. It then compares the
entire generated set with the coordinate root set, in seed coordinates.
We check $2\leq k\leq8$ in types $\typ B,\typ C$, including
$\typ C_2$, and all of $\typ F_4$. The numbers of generated roots
are $2k^2$, $2k^2$ and 48, respectively; their numbers of lines
are $k^2$, $k^2$ and 24. In rank 2 the angle $60^\circ$
does not occur in the $\typ B,\typ C$ line set.

\subsection{Containing and noncontaining examples}

The following integer Gram matrices give small sets satisfying the
partition condition. The diagonal specifies the assigned squared
lengths. Each count is the number of roots after reflection closure;
the initial set consists of one vector per matrix row.
\begin{center}
\begin{tabular}{@{}lcc@{}}
\toprule
Example&Seed Gram matrix&Roots in the closure\\\midrule
Empty set&empty&0\\[2pt]
One line&$(2)$ or $(4)$&2\\[2pt]
Two orthogonal lines&$\opk{diag}(2,4)$&4\\[2pt]
Two $60^\circ$ lines&$\left(\begin{smallmatrix}2&-1\\-1&2\end{smallmatrix}\right)$&6\\[3pt]
Two $45^\circ$ lines&$\left(\begin{smallmatrix}2&2\\2&4\end{smallmatrix}\right)$&8\\[3pt]
Three lines&$\left(\begin{smallmatrix}2&2&1\\2&4&2\\1&2&2\end{smallmatrix}\right)$&18\\[3pt]
Three dependent lines&$\left(\begin{smallmatrix}2&-1&-1\\-1&2&-1\\-1&-1&2\end{smallmatrix}\right)$&6\\
\bottomrule
\end{tabular}
\end{center}
For instance, the three-line matrix with 18 roots is realized by
$\Ze_1+\Ze_2$, $2\Ze_1$, $\Ze_1+\Ze_3$ in $\re^3$ and
generates $\typ C_3$. The last matrix is realized by
$\Ze_1-\Ze_2$, $\Ze_2-\Ze_3$, $\Ze_3-\Ze_1$; its span is a
plane. These include sets which are not initially reflection-closed.
The checks also place a two-line $45^\circ$ configuration in a plane
in $\re^4$ and add two orthogonal seeds of squared length 2, as in
Lemma~\ref{lem:seed-completion}. Reflection closure has 12 roots,
while $\Psi(T')$ has 32; the script asserts rank 4 for both outputs.
A further check uses
$\typ B_2\oplus3\typ B_2$. Its four squared lengths are
$1,2,9,18$, so this example checks independent component scales.

For noncontainment, put $\gamma=1/\sqrt2$ and consider the two unit Gram
matrices
\[
 M_1=\pmat{1&\gamma&\gamma\\\gamma&1&\gamma\\\gamma&\gamma&1},\qquad
 M_2=\pmat{1&1/2&\gamma\\1/2&1&1/2\\\gamma&1/2&1}.
\]
Both are realized by three independent real unit vectors.
For $M_1$, explicit vectors are
\[
 \Ze_1,\quad \gamma\Ze_1+\gamma\Ze_2,\quad
 \gamma\Ze_1+(1-\gamma)\Ze_2+\sqrt{2\gamma-1}\Ze_3.
\]
Here $2\gamma-1>0$; their squared norms are 1, and the three inner
products are $\gamma,\gamma,\gamma^2+\gamma(1-\gamma)=\gamma$.
For $M_2$, use
\[
 \Ze_1,\quad \tfrac12\Ze_1+\tfrac{\sqrt3}{2}\Ze_2,\quad
 \gamma\Ze_1+\tfrac{1-\gamma}{\sqrt3}\Ze_2
       +\sqrt{\tfrac12-\tfrac{(1-\gamma)^2}{3}}\Ze_3.
\]
The last radicand is positive because $0<1-\gamma<1$. The third vector
has squared norm $\gamma^2+1/2=1$, and its inner products with the first
two are $\gamma$ and $\gamma/2+(1-\gamma)/2=1/2$.
Thus the matrices describe actual line sets, not just proposed angles.

In $M_1$, every pair has angle $45^\circ$. A compatible partition
would put every pair in different parts, which is impossible on three
lines. In $M_2$, the two $60^\circ$ pairs force all three lines into
the same part, while the $45^\circ$ pair forces two of them into
different parts. Thus neither configuration is contained in a Root
System, by the necessary direction. The same obstruction applies in
any larger ambient space. The second example also shows why testing
only bipartiteness of the graph of $45^\circ$ pairs is insufficient:
its graph has just one edge. One must also enforce equality of parts
along $60^\circ$ pairs.

\subsection{What the computations establish}

Run the checks from the repository root with
\begin{quote}
\ml{python cgss/proofs/thm312/checks/check_thm312.py},
\end{quote}
for a Python 3 with SymPy and NumPy. The script prints its results and
writes no files. We ran it with Python 3.12.14, SymPy 1.14.0 and NumPy;
its output records the installed versions. Integers, rational numbers and exact
algebraic numbers determine every acceptance decision. Numerical
eigenvalues are additional diagnostics for the two noncontaining
examples, not the criterion for positive definiteness.

The closure calculation stores integer coefficient vectors $a$ in the
seeds and uses $Ma$ as the key for a represented vector. This also
handles dependent seeds: if $M$ is their Gram matrix, then
$Ma=Mb$ says that the represented difference is orthogonal to every
seed. It lies in their span, so it is zero. Thus different coefficient
expressions of one vector are identified. For each resulting set the
script checks nonzero vectors, both signs, spanning, reducedness,
integrality, reflection closure and squared lengths. For the normalized
seeds it also checks even inner products of long roots with every seed.
Evenness against all of $L$ follows by integer linearity, as in the proof.

The explicit set $\Psi(T)$ is checked separately for the independent
seed sets of $\typ A_1^2$, $\typ A_1^3$, $\typ A_2\oplus\typ A_1$,
$\typ B_2$, $\typ C_3$, $\typ D_4$, the ambient-completion example,
and the empty set. To enumerate it completely, for a positive definite
seed Gram matrix $M$ and an integer coefficient vector $a$, the script
uses
\[
 q\left(\sum_i a(i)\Zt_i\right)\leq4
 \quad\Longrightarrow\quad a(i)^2\leq4M^{-1}(i,i).
\]
This is Cauchy--Schwarz against the dual basis. The resulting finite
coefficient box contains every candidate, so enumeration has a proved
bound. The script tests the parity condition and all Root System axioms,
and compares the result with reflection closure.

There are exactly $5^3=125$ labelled symmetric $3\times3$ matrices
with diagonal 1 and each upper off-diagonal entry in
$\{0,1/2,-1/2,1/\sqrt2,-1/\sqrt2\}$. The script checks all of them.
Their proper principal minors are positive; the exact determinant
decides positive semidefiniteness. The 97 feasible matrices have the
following cross-tabulation, which the script asserts and prints:
\begin{center}
\begin{tabular}{@{}lrrr@{}}
\toprule
&With partition&Without partition&Total\\\midrule
Positive definite&65&16&81\\
Singular&16&0&16\\\midrule
Total&81&16&97\\\bottomrule
\end{tabular}
\end{center}
A graph propagation algorithm and exhaustive enumeration
of all eight binary colourings agree in every case. For each accepted
configuration the script constructs and verifies the reflection closure
for both assignments of 2 and 4 to its parts. The 162 resulting closures
have sizes $6,8,10,12,18$, occurring $10,36,12,32,72$ times, respectively.
Independently of those colourings, the script tries all 27 assignments
of squared lengths in $\{1,2,4\}$ for each of the 97 feasible unit
Gram matrices. It checks both ordered crystallographic coefficients
for every pair of seeds. Each passing assignment is closed under
reflections without assuming the normalized parity condition, and
the output is checked against every Root System axiom, with its rank
asserted. Some integral assignments have nonreduced closures; these are
rejected as Root System candidates. Amongst the 2619 assignments, 2190
fail seed integrality, 24 give nonreduced closures, and 405 give Root
Systems. In every configuration, an integral assignment exists if
and only if a compatible partition exists, and if and only if an
assignment has a verified Root System completion. In particular, all
27 assignments fail seed integrality for each of the 16 configurations
without a partition. These are finite obstruction checks independent
of the colouring calculation. The script also checks the parity-failure
example above.

With random seed 20260928, the script samples 3000 proper subsets of
2 to 9 lines from each of
\[
 \typ F_4\oplus\sqrt3\,\typ B_3,
 \qquad \typ C_3\oplus\typ G_2\oplus\tfrac{17}{10}\typ B_2.
\]
It verifies the full containing Root Systems first. The $\typ G_2$
model is the one in \cite[coordinate-model section]{R311}, in the
plane of $\re^3$ with coordinate sum zero. Rational coordinates with
block inner-product weights 3 and $289/100$ represent the two rescalings
exactly. The 1014 subsets with a $30^\circ$ pair are counted and skipped;
all 4986 remaining subsets are checked by graph propagation and by the parity of
the integer logarithm of their actual squared root lengths. Common
rescaling of squared lengths by $2^{-4}$ and $2^{-3}$ also checks negative
logarithms and both parity shifts. The output records the counts and
asserts coverage of $45^\circ$ pairs, multiple components and both
length classes of the $\typ G_2$ component.

These computations check finite instances, including obstructions,
dependent seeds and degenerate cases. They neither prove the
unbounded-rank theorem nor supply Lean proof terms. Finiteness in all
dimensions is proved in Lemma~\ref{lem:integer-span} and
Proposition~\ref{prop:psi}. For the normalized seeds, reflection
closure stays inside the finite reflection-closed set $\Psi(T)$,
which also proves termination of that comparison calculation.
The ordinary proof does not use any of the computed counts.

\section{Sources and the mathematical interface}\label{sec:sources}

We will record the precise division between existing facts and the
arguments proved here. Our Mathlib reference \cite{Mathlib} is Mathlib at commit
\begin{center}
\ml{082e2d37e8b0463410cdb532e111cd43d5a66174}
\end{center}
(4 September 2026), the version pinned in \ml{lake-manifest.json}.
The names below were checked against its source files.

\begingroup\raggedright
\begin{enumerate}
\item From \cite{CGSS} we use the meaning of a Root System in
 Definition 3.9 and the statement to be proved in Theorem 3.12,
 corrected as in Section~\ref{sec:statement}. No theorem of
 \cite{CGSS}, including Theorem 3.10, is required as a premise of
 our proof. In particular, we do not assume a completion theorem.
\item From \cite{R311}, Theorem C supplies the pairwise root-angle
 assertion; Lemma~\ref{lem:root-angles} repeats the needed proof
 from the integrality axiom. Its Theorem B and its verified coordinate
 models identify the $\typ B,\typ C,\typ F$ systems used for checks.
 Theorems A and B are not premises of the equivalence. Thus no
 root-system classification, or two-length theorem for irreducible
 systems, is needed in the main argument.
\item Finite-dimensional linear algebra supplies a basis chosen
 from a spanning set. In
 \ml{Mathlib/LinearAlgebra/LinearIndependent/Lemmas.lean},
 \ml{exists_linearIndependent} gives an independent subset with the
 same span over a division ring. The declaration
 \ml{LinearIndependent.finite} in
 \ml{Mathlib/LinearAlgebra/Dimension/Finite.lean} makes that subset
 finite in a finite-dimensional space. We apply this over $\re$ in the two finiteness
 lemmas. The empty span gives the empty basis.
 For the line interface, \ml{finrank_eq_one_iff'} in
 \ml{Mathlib/LinearAlgebra/Dimension/FreeAndStrongRankCondition.lean},
 applied to the subtype $\ell$, supplies a nonzero vector on a
 one-dimensional subspace. The declaration
 \ml{eq_span_singleton_of_mem_of_finrank_eq_one} in
 \ml{Mathlib/LinearAlgebra/FiniteDimensional/Basic.lean} says that
 $\ell=\opk{span}_{\re}\{\Zr\}$ for any nonzero $\Zr\in\ell$.
 In the same file, \ml{finrank_span_singleton} proves the converse
 dimension assertion for a nonzero vector. These facts justify the
 chosen representatives and the definition of \ml{rootLines} below.
\item Bilinearity, positive definiteness and Cauchy--Schwarz are in
 \ml{Mathlib/Analysis/InnerProductSpace/Basic.lean}.
 The declarations \ml{real_inner_self_eq_norm_sq},
 \ml{inner_self_eq_zero} and \ml{abs_real_inner_le_norm} apply to
 real inner product spaces and give exactly the norm facts used
 above. Reflection identities are expanded in
 Lemma~\ref{lem:reflection}, rather than taken from an API with a
 possibly different reflection convention.
\item Integer-span induction is supplied by
 \ml{Submodule.span_induction} and
 \ml{Submodule.span_induction}\textsubscript{2}
 in \ml{Mathlib/LinearAlgebra/Span/Defs.lean}, with scalars $\ints$.
 We use their ordinary content, finite integer linear combinations,
 in Lemma~\ref{lem:integer-span}. Finiteness from an injection into
 a finite set is \ml{Set.Finite.of_injOn} in
 \ml{Mathlib/Data/Set/Finite/Basic.lean}. These facts justify both
 finite-target arguments, without a lattice-discreteness theorem.
\item The irrationality used for reducedness is
 \ml{irrational_sqrt_two} in
 \ml{Mathlib/NumberTheory/Real/Irrational.lean}. Strict increase of
 integer powers of $2\in\re$ is
 \ml{zpow_right_strictMono}\textsubscript{0} in
 \ml{Mathlib/Algebra/Order/GroupWithZero/Basic.lean}, with hypothesis
 $1<2$. In \ml{Mathlib/Data/Int/Log.lean},
 \ml{Int.zpow_log_le_self} and \ml{Int.lt_zpow_succ_log_self}
 give \eqref{eq:log-bounds} for \ml{Int.log 2 x} when $x>0$.
 The equivalent bounds \ml{Int.zpow_le_iff_le_log} and
 \ml{Int.lt_zpow_iff_log_lt}, applied at $2x>0$, turn the multiplied
 inequalities into $\eta(x)+1\leq \eta(2x)<\eta(x)+2$.
 Since these are integers, this proves \eqref{eq:log-shift} in Lean
 without an analytic-logarithm API.
\item An orthonormal basis of a finite-dimensional inner product
 space is supplied by \ml{exists_orthonormalBasis}, or the indexed
 choice \ml{stdOrthonormalBasis}, in
 \ml{Mathlib/Analysis/InnerProductSpace/PiL2.lean}.
 We apply it to $W^\perp$ with its induced inner product.
 Orthogonal decomposition is \ml{Submodule.isCompl_orthogonal}
 in \ml{Mathlib/Analysis/InnerProductSpace/Projection/Basic.lean}.
 It assumes \ml{HasOrthogonalProjection}; the instance
 \ml{Submodule.HasOrthogonalProjection.ofCompleteSpace} supplies
 this for $W$, since a finite-dimensional real normed space is
 complete. Thus its hypotheses are satisfied
 even when $W=0$ or $W=V$.
\end{enumerate}
\par\endgroup

All remaining arithmetic consists of
finite sums, the positive integer factor pairs of 1 and 2, or parity
of integers. Thus Theorem~\ref{thm:main} rests only on the
listed elementary library facts and the explicit proofs in this document.

\section{Lean statement shape}\label{sec:lean}

We will keep the settled name
\ml{LineSystem.subset_rootSet_iff_bipartition} from
\ml{cgss/notation.md}, Section 10. There is a necessary distinction
between the namespace and the argument type. The current
\ml{LineSystem V}, defined in
\ml{LineGraphs/LineSystem/Defs.lean}, already requires all normalized
inner products to be in $\{-2,-1,0,1,2\}$. It represents only
$60^\circ$ and $90^\circ$ lines. Using an argument of that type
would omit all $45^\circ$ cases of this theorem.

We will state the theorem in the \ml{LineSystem} namespace
on an unbundled set
\[
 S:\ml{Set (Submodule Real V)},
\]
with \ml{[NormedAddCommGroup V]}, \ml{[InnerProductSpace Real V]}
and \ml{[FiniteDimensional Real V]}. We will require
\[
 \ml{hline}:\quad
 \forall\ell\in S,\quad\ml{Module.finrank}\ \re\ \ell=1.
\]
The following auxiliary predicates give a precise statement without
choosing representatives in its type. They are definitions of the
project, in \ml{LineGraphs/Theorem312/Defs.lean}, not declarations of Mathlib.

For subspaces $\ell,m$ and $a\in\re$, the predicate
\ml{LineSystem.HasSquaredCosine} $\ell\ m\ a$ means
\[
 \forall\Zx\in\ell\ \forall\Zy\in m,\qquad
 (\Zx,\Zy)^2=a\,q(\Zx)q(\Zy).
\]
For lines this is equivalent to \eqref{eq:cosine} having value $a$:
nonzero representatives give the forward implication, and multiplying
representatives gives the reverse implication, including zero vectors.
The division-free form avoids nonzero proofs in the predicate itself.
Define \ml{LineSystem.HasAllowedAngles} $S$ to mean
\[
 \forall\ell,m\in S,\ \ell\ne m\ \Longrightarrow
 \exists a\in\{0,1/4,1/2\},\quad
       \ml{HasSquaredCosine}\ \ell\ m\ a.
\]
Define \ml{LineSystem.HasBipartition} $S$ to mean that there is a
function $b:S\to\ml{Bool}$ such that for distinct $\ell,m\in S$,
\[
\begin{cases}
 \ml{HasSquaredCosine}\ \ell\ m\ 0\ \lor
 \ml{HasSquaredCosine}\ \ell\ m\ (1/4),&b(\ell)=b(m);\\
 \ml{HasSquaredCosine}\ \ell\ m\ 0\ \lor
 \ml{HasSquaredCosine}\ \ell\ m\ (1/2),&b(\ell)\ne b(m).
\end{cases}
\]
The two fibres of $b$ are exactly the parts of condition (2), and
neither fibre is required to be nonempty.

The predicate \ml{IsRootSet} on $\Phi:\ml{Set V}$ contains
the six requirements of Definition~\ref{def:rootset}: finiteness,
nonzero vectors, real span equal to the top submodule,
crystallographic integrality, reducedness and reflection closure.
Define \ml{LineSystem.rootLines} $\Phi$ by
\[
 \{\ell:\exists\Zr\in\Phi,\quad
              \ell=\opk{Submodule.span}_{\re}\{\Zr\}\}.
\]
The nonzero-root field ensures that these are lines whenever
\ml{IsRootSet} $\Phi$ holds. In particular the statement is line
containment, not containment of a particular normalization of root vectors.

With these definitions the declaration has the following conclusion,
under the displayed typeclasses and hypotheses \ml{hline} and
\ml{hangle : HasAllowedAngles S}:
\begin{equation*}
 \bigl(\exists\Phi:\ml{Set V},\quad
     \ml{IsRootSet}\ \Phi\ \land\ S\sbs\ml{rootLines}\ \Phi\bigr)
 \quad\Longleftrightarrow\quad \ml{HasBipartition}\ S.
\end{equation*}
The theorem's full name is
\ml{LineSystem.subset_rootSet_iff_bipartition}.
Finiteness of $S$ is derived from \ml{hline} and \ml{hangle}, as
in Lemma~\ref{lem:finite-lines}; a finite-set wrapper can expose it
explicitly. Spanning and nonemptiness of $S$ are not hypotheses.
The docstring of the theorem records these conventions and the errata
pointers:
\begin{quote}\small
Theorem 3.12 of [CGSS]. For one-dimensional subspaces with squared
cosines 0, 1/4 or 1/2, root-line containment is equivalent to a
compatible bipartition. The set and either colour fibre may be empty
(Implicit hypothesis 8); finiteness is derived and no spanning
hypothesis is imposed on the lines. The containing root set spans the
given ambient space. The cross-part condition reads ``or'' (Error 13,
item 3); the explicit construction fills Gap 11.
\end{quote}

For the construction, use \ml{Submodule.span Int T} for $L$ and
define \ml{LineSystem.integralSpanRoots} $T$ by the set-builder
\eqref{eq:psi}. The squared-length-4 condition quantifies over every
element of this fixed $L$. Use the ambient version of
Proposition~\ref{prop:psi}, with the extra hypothesis
$\opk{span}_{\re}(T)=V$. First form $T'$ by
Lemma~\ref{lem:seed-completion}, using \ml{stdOrthonormalBasis}
on $W^\perp$ and viewing its vectors in $V$. Then apply the proposition
to $T'$ directly. This proof uses no \ml{IsRootSet} on the subtype
$W$ and no extension lemma for a union of Root Systems.
The span assertion of Corollary~\ref{cor:normalized-containment}
can be formalized separately by restricting the same construction to $W$.
The helper theorems in the \ml{LineSystem} namespace are
\ml{finite_of_hasAllowedAngles}, \ml{isRootSet_integralSpanRoots}
and \ml{exists_rootSet_of_hasBipartition}, following the conclusion
rule in the notation file.

For the necessary direction, set the Boolean colour to
\ml{decide (Int.emod (Int.log 2 (inner r r)) 2 = 1)} for the chosen root
on a line. Equations \eqref{eq:log-bounds} and \eqref{eq:log-shift}
give the argument without components or path choices.

The namespace follows rule 2 of \ml{cgss/notation.md}, Section 10,
which places line-system results in \ml{LineSystem}; the argument is
the unbundled set $S$, not a term of the structure \ml{LineSystem V}.
The notation file records this shape and the five definitions
\ml{HasSquaredCosine}, \ml{HasAllowedAngles}, \ml{HasBipartition},
\ml{rootLines} and \ml{integralSpanRoots}.

For a bridge to Section 2, if $R:\ml{LineSystem V}$, then
\ml{rootLines} $R.\Phi$ satisfies \ml{HasAllowedAngles} and
\ml{HasBipartition}, with the constant Boolean colouring.
Indeed, its vectors have squared length 2; inner products $\pm2$
give equal lines, since the squared norm of their difference or sum
is zero. Distinct lines thus have inner product $0$ or $\pm1$ on
these representatives, giving squared cosine $0$ or $1/4$.
Equation~\eqref{eq:cosine} transfers this to all representatives.
The bridge lemmas are \ml{LineSystem.hasAllowedAngles_rootLines}
and \ml{LineSystem.hasBipartition_rootLines}.

No bridge to \ml{RootPairing} is needed to prove this statement.
If the bridge proved in \cite{R311} is used afterwards, Mathlib's pairing has the
squared length of the \emph{second} root in the denominator:
\[
 p(\Zr,\Zs)=\frac{2(\Zr,\Zs)}{q(\Zs)}.
\]
Definition~\ref{def:rootset} uses the first root as the reflecting
root, so the bridge interchanges the arguments. Integrality for all
ordered pairs provides both conventions. The Root System produced
here is reduced and crystallographic and spans its stated ambient
type, exactly as required by the bridge in \cite{R311}.

There is no unresolved mathematical step in the ordinary proof.
The structured version of these proofs is
\ml{cgss/proofs/thm312/thm312-lamport.tex}, and the Lean formalization, in
\ml{LineGraphs/Theorem312/}, proves
\ml{LineSystem.subset_rootSet_iff_bipartition} with the helper
definitions and theorems above; it uses no computational check as a
mathematical premise.

\begin{thebibliography}{Mathlib}
\raggedright
\bibitem[CGSS]{CGSS}
P. J. Cameron, J. M. Goethals, J. J. Seidel and E. E. Shult,
\emph{Line graphs, root systems, and elliptic geometry},
J. Algebra \textbf{43} (1976), 305--327.
Definition 3.9 and Theorem 3.12, pp.~314--315; see also
\ml{cgss/errata.md}, Error 13, item 3 and Gap 11.

\bibitem[R311]{R311}
\emph{The classification behind Remark 3.11 of [CGSS]},
project manuscript \ml{cgss/proofs/remark311/remark311.tex},
Theorems A--C and the coordinate-model and Mathlib-interface sections.

\bibitem[Mathlib]{Mathlib}
The mathlib community, \emph{Mathlib}, commit
\ml{082e2d37e8b0463410cdb532e111cd43d5a66174} (4 September 2026),
the version pinned in \ml{lake-manifest.json}.
The specific source files and declarations used here are listed in
Section~\ref{sec:sources}.
\end{thebibliography}
\end{document}
